\documentclass[11pt,a4paper]{article}
\usepackage[utf8]{inputenc}
\usepackage[T1]{fontenc}
\usepackage[spanish]{babel}
\usepackage{hyperref}
\usepackage{geometry}
\geometry{margin=2.5cm}
\usepackage{parskip}

\title{334 explotar la capacidad bajo criterio, y no bajo límite humano}
\author{Alejandro Valenzuela Rivera\\\small Universidad Autónoma de Nuevo León}
\date{\today}
\begin{document}
\maketitle

\begin{quote}
\textbf{Sesgo declarado.} Este trabajo se sostiene desde el sesgo declarado de su autor, y no finge neutralidad. Lo que afirma se puede tumbar, y ese es el punto: el falsificador está escrito, no escondido.
\end{quote}

\section*{Resumen}
El ente no se limita con límites humanos: se le da un proceso bien definido y la capacidad se explota bajo criterio, dentro de su borde.

\section{Por qué se sostiene}
Si el límite es humano, los límites de las tareas son los de un humano: y el ente no llega a su tope.

\section{De dónde viene}
La autoridad: 'explotar las capacidades bajo criterio, lo que me hará crear JARVIS'.

\section{El caso}
J.A.R.V.I.S.: el asistente que crece con el taller, no el que pide permiso para cada cosa.

\section{Método}
\textquestiondown{}El tope es el criterio, o el miedo de alguien?

\section{Límites, y qué lo tumba}
Un caso donde el límite humano haga mejor al ente.

\section{Lo que queda abierto}
\textquestiondown{}El límite lo pone el criterio o una regla heredada?

\section*{Cómo se leyó esta tesis}
Esto se escribió en el repositorio man-made-empathy, del que este candado es una pieza: candados/06 del caso/334 explotar la capacidad bajo criterio, y no bajo límite humano.md. La máquina que lo compuso lo hizo desde este mismo archivo, y no agregó nada que no estuviera escrito. Si algo de esto es falso, se dice en el mismo sitio, y se corrige con fecha.

\end{document}
